\documentclass[10pt,a4paper]{amsart}
\usepackage[margin=19mm]{geometry}
\usepackage{amsmath,amssymb,mathtools}
\usepackage[scr=rsfs,cal=euler]{mathalfa}
\usepackage{xfrac}
\usepackage[unicode,hidelinks,bookmarks=false]{hyperref}
\newcommand{\ie}{i.e.\ }
\newcommand{\cf}{cf.\ }
\newcommand{\resp}{respectively\ }
\newcommand{\Subsection}{\subsection}
\providecommand{\nobreakdash}{\nobreak\hbox{-}}
\setlength{\emergencystretch}{2em}
\multlinegap0pt
\usepackage{bm}
\usepackage{bbm}
\usepackage{dsfont}
\usepackage{xspace}
\usepackage{cancel}
\usepackage{mathtools}
\usepackage{amsthm}
\makeatletter
\@ifundefined{theorem}{\newtheorem{theorem}{Theorem}}{}
\@ifundefined{lemma}{\newtheorem{lemma}[theorem]{Lemma}}{}
\@ifundefined{proposition}{\newtheorem{proposition}[theorem]{Proposition}}{}
\makeatother
\def\F{\mathbb{F}} 
\title[Comaximal Graphs of finite-dimensional Lie algebras over fin]{Comaximal Graphs of finite-dimensional Lie algebras over finite fields: Triangle counts and structural invariants}
\author{David Towers, Yesneri Zuleta, Ismael Gutierrez}
\date{Source: arXiv:2608.16575v1 (CC BY 4.0)}
\begin{document}
\maketitle

\noindent\emph{Source and licence.} Comaximal Graphs of finite-dimensional Lie algebras over finite fields: Triangle counts and structural invariants. Version arXiv:2608.16575v1 (CC BY 4.0), distributed under the Creative Commons Attribution 4.0 International licence, as recorded in the arXiv licence metadata for this exact version and shown on the abstract page https://arxiv.org/abs/2608.16575v1. Full rights notice: licenses/arxiv-2608.16575-CC-BY-4.0.txt.\par\smallskip

\noindent\textit{Editorial scope.} Counted unit: the Proposition whose source label is \texttt{B2} in the article (source0.tex lines 360--369, source label \texttt{B2}), delivered together with the article's own complete original proof of it (source0.tex lines 371--470, one \texttt{proof} environment of 100 lines). No mathematics is written, repaired or re-derived: statement and proof are the article's own text line for line. Required context retained uncounted: T3P (lines 240--242). Every definition, display and label of those lines is retained unchanged. Omitted from this package and not redistributed: 1--239, 243--359, 370--370, 471--1458. Nothing inside a retained excerpt is omitted or altered: every retained body is delivered exactly as the article's lines are written, so no omission receipt of any kind accompanies this package. Citations: the retained excerpts carry no citation command at all, so no bibliography entry of the article is transcribed and no reference is redistributed. Reference adaptation: none was needed and none was made; every reference inside the delivered unit resolves inside it, so no unresolved reference or citation is produced. Printed slips kept literal: no printed word, symbol, hypothesis or number of the counted statement or of its proof was altered. Layout and class: the article's own class is \texttt{amsart} with packages including amsmath, amsthm, amscd, amssymb, amsfonts, color, graphicx, mathrsfs, figsize, algorithm, algpseudocode, enumerate, float, tikz; the wrapper is the lane's pinned \texttt{amsart} editorial wrapper at 10pt on A4, which restates the theorem-like environments the retained text needs and adds the article's own macro definitions that the retained text needs (\texttt{\textbackslash F}), with extra wrapper packages (bm, bbm, dsfont, xspace, cancel, mathtools). The delivered numbering is the unit's own. Assets: no retained span contains a figure environment, a graphics-inclusion command, a PDF-page inclusion, a table or a TikZ picture; no image file is copied into this package, the package declares no asset dependency and no image file is redistributed with it. Licence: version arXiv:2608.16575v1 of this article is distributed under the Creative Commons Attribution 4.0 International licence, as recorded in the arXiv licence metadata for this exact version and shown on the abstract page https://arxiv.org/abs/2608.16575v1. A full rights notice is delivered at licenses/arxiv-2608.16575-CC-BY-4.0.txt. Characters: every retained excerpt is pure ASCII, so the delivered engine, XeLaTeX with its default Latin Modern text font, reports no missing character.
\begin{lemma}\label{T3P}
Let $\Gamma$ be a graph in which the induced subgraph on a set $\mathcal{P}$ of vertices is a complete graph $K_{m}$.  Then the number of triangles of $\Gamma(L)$ whose three vertices all lie in $\mathcal{P}$ is $\tbinom{m}{3}$.
\end{lemma}

\begin{proposition}\label{B2}
Let $L$ be the three-dimensional solvable non-nilpotent Lie algebra over $\F_q$ with $\dim L'=1$. Then
\[\mathfrak{t}(\Gamma(L)) = \frac{q(q+1)(q^4+2q^3+5q^2-6q-2)}{6}.\]
% \begin{align*}
% \mathfrak{t}(\Gamma(L)) =
% &   \frac{q(4q^2-1)}{3} + q^2(q-1)(2q+1) + \frac{q^2(q-1)^2(2q+1)}{2} + \frac{q^2(q-1)^3}{2} \\
% &+ \frac{q^2(q-1)^2(q-2)^2}{6}.
% \end{align*}

\end{proposition}

\begin{proof}
We classify triangles by vertex type.

\medskip
\noindent\textbf{1. Three planes.} By Lemma \ref{T3P} with $m=2q+1$ we have: $T_{3P}= \binom{2q+1}{3} =\frac{q(4q^2-1)}{3}$.

\smallskip
\noindent\textbf{2. Two planes and one line.}
%Since $\langle y\rangle$ and $\langle z\rangle$ are adjacent to no line, we separate into seven pair types.
For a line $A$ and a plane $P$ of $L$,  the structural description is as follows:

\begin{enumerate}
\item[(1)] The lines $\langle y \rangle$ and $\langle z \rangle$ are adjacent
to no other line, but $\langle y \rangle \sim P_\alpha$ for every
$\alpha \in \mathbb F_q$, and $\langle z \rangle \sim Q_b$ for every
$b \in \mathbb F_q$; neither is adjacent to $V$, and $\langle y \rangle$ is
not adjacent to any $Q_b$, nor $\langle z \rangle$ to any $P_\alpha$.
\item[(2)] The $q-1$ intermediate lines $I_\alpha = \langle \alpha y + z
\rangle$, $\alpha \in \mathbb F_q^*$, each adjacent to every line outside
$V$ but to no line inside $V$.
\item[(3)] The $q^2$ lines $X_{a,b} = \langle x + ay + bz \rangle$,
$a,b \in \mathbb F_q$, where $X_{a,b} \sim X_{c,d}$ if and only if $a \neq c$
and $b \neq d$.
\end{enumerate}

We separate into five pair types:
\begin{enumerate}
\item[(i)] Pair $(Q_{b_1},\,Q_{b_2})$. The $q-1$ intermediates lie outside both planes; each $X_{a,b}$ satisfies $X_{a,b}\sim Q_{b_i}$ if and only if $b\neq b_i$, so
$b\notin\{b_1,b_2\}$ gives $q(q-2)$ valid $X$-lines.  In addition, $\langle z \rangle$ is adjacent to
every $Q_b$, hence to both $Q_{b_1}$ and $Q_{b_2}$. Including this line, the count per pair is $q^2-q$. With
$\binom{q}{2}$ pairs: $T_{2PL}^{(i)} = \binom{q}{2}\,q(q-1)$.

% The $q-1$ intermediate lines are not contained in any $Q_b$ plane, so each is adjacent to both $Q_{b_i}$. An $X$-line $X_{a,b}$ is adjacent to $Q_{b_i}$ if and only if $b\neq b_i$, so it is adjacent to both if and only if $b\notin\{b_1,b_2\}$, giving $q(q-2)$ such lines. The count per pair is $(q-1)+q(q-2)=q^2-q-1$, hence
% \[T_{2PL}^{(i)}=\binom{q}{2}(q^2-q-1) = \frac{1}{2}q(q-1)(q^2-q-1).\]

\smallskip

\item[(ii)] Pair $(P_{\alpha_1},\,P_{\alpha_2})$. By the symmetric argument
(swap $a$- and $b$-coordinates), together with the contribution of
$\langle y \rangle$ (adjacent to every $P_\alpha$), the count per pair is
also $q^2-q$, giving $T_{2PL}^{(ii)} = T_{2PL}^{(i)}$.

% All the $q-1$ intermediate lines are adjacent to both planes. An $X$-line $X_{a,b}$ is adjacent to $P_{\alpha_i}$ if and only if $a\neq\alpha_i$, so it is adjacent to both if and only if $a\notin\{\alpha_1,\alpha_2\}$, giving $q(q-2)$ such lines. The count per pair is again $q^2-q-1$, hence $T_{2PL}^{(ii)}= T_{2PL}^{(i)}$. 
\smallskip

\item[(iii)] Pair $(Q_b,\,P_\alpha)$. In this case, their intersection is the line $X_{\alpha, b}$, and all $q-1$ intermediate lines are adjacent to both.
An $X$-line $X_{a,b'}$ is simultaneously adjacent to $Q_b$, and to $P_\alpha$
if and only if $b'\neq b$ and $a\neq\alpha$, giving $(q-1)^2$ such lines, and
neither $\langle y \rangle$ nor $\langle z \rangle$ is adjacent to both $Q_b$
and $P_\alpha$, since $\langle y \rangle \not\sim Q_b$ and $\langle z \rangle
\not\sim P_\alpha$). Hence $T_{2PL}^{(iii)} = q^3(q-1)$.  
\smallskip

\item[(iv)] Pair $(V, Q_b)$. Only $X$-lines are adjacent to $V$; $\langle z
\rangle$ is not adjacent to $V$, giving $q(q-1)$ lines adjacent to both. Thus, $T_{2PL}^{(iv)}= q^2(q-1)$.
\smallskip

\item[(v)] Pair $(V, P_\alpha)$. Symmetric to (iv) we have  $T_{2PL}^{(v)}= q^2(q-1)$.
\end{enumerate}
Summing the five terms, we have
\[T_{2PL} = q^2(q-1)^2 + q^3(q-1) + 2q^2(q-1) = q^2(q-1) ((q-1)+q+2) = q^2(q-1)(2q+1).\]

\noindent\textbf{3. One plane and two lines.} We compute the number of edges among lines adjacent to a plane $P$, for each plane type. We denote it by $e_P$. Then we have $T_{P2L}=\sum_{P\in \mathcal{P}} e_P$. Fix a plane $P\in \mathcal{P}$.
\begin{enumerate}
\item[(i)] $P=V$. Lines adjacent to $V$ are the $q^2$ $X$-lines.  Non-adjacent pairs share an $a$-value ($q\binom{q}{2}$ pairs) or a $b$-value (another $q\binom{q}{2}$ pairs), with the two events disjoint.  Hence
\begin{equation}\label{eV-B2}
  e_{V}=\binom{q^2}{2}-2q\binom{q}{2}=\frac{q^2(q-1)^2}{2}.
\end{equation}

\smallskip

\item[(ii)] $P= Q_b$. Lines adjacent to $Q_{b}$: $q-1$ intermediates and $q(q-1)$ $X$-lines with $b'\neq b$.  No two intermediates are adjacent ($0$ edges); each intermediate is adjacent to every $X$-line ($q(q-1)^2$ edges); among $X$-lines, non-adjacent pairs share an $a$-value ($q\binom{q-1}{2}$) or a $b'$-value ($(q-1)\binom{q}{2}$), giving
\begin{equation}\label{e2X-B2}
e_{2X} =\binom{q(q-1)}{2}-q\binom{q-1}{2}-(q-1)\binom{q}{2} =\frac{q(q-1)^2(q-2)}{2}.
\end{equation}
Hence $e_{Q_{b}}=q(q-1)^2+\frac{q(q-1)^2(q-2)}{2}=\frac{q^2(q-1)^2}{2}$.

\smallskip

\item[(iii)] $P= P_{\alpha}$. By the symmetric argument (swapping $a$- and $b$-coordinates) follows $e_{P_{\alpha}} =\frac{q^2(q-1)^2}{2}$.
\end{enumerate}
Every plane gives $e_{P}=\frac{q^2(q-1)^2}{2}$, so
\[T_{P2L}=(2q+1) \frac{q^2(q-1)^2}{2} =\frac{q^2(q-1)^2(2q+1)}{2}.\]
Note that neither $\langle y\rangle$ nor $\langle z \rangle$ is adjacent to any line, their presence
among the vertices adjacent to a fixed plane $P$ contributes no new edges. 

\smallskip
\noindent\textbf{4. Three lines.} Since $\langle y \rangle$ and $\langle z
\rangle$ are adjacent to no line, $\langle y\rangle$, $\langle z\rangle$ are in no triangle; no two intermediates are adjacent. Then, in this case, only two options arise.
\begin{enumerate}
\item[(i)] One intermediate line and two $X$-lines. Any adjacent $X$-pair forms a triangle with each of the $q-1$ intermediates, giving
\[T_{3L}^{(i)} = (q-1) \frac{q^2(q-1)^2}{2}=\frac{q^2(q-1)^3}{2}.\]

\item[(ii)] Three $X$-lines. $X_{a,b},X_{c,d},X_{e,f}$ are mutually adjacent if and only if $a,c,e$ are pairwise distinct and $b,d,f$ are pairwise distinct.  Counting ordered triples and dividing by $3!$: 
\[T_{3L}^{(ii)} =\frac{(q(q-1)(q-2))^2}{6}=\frac{q^2(q-1)^2(q-2)^2}{6}.\]
\end{enumerate}
Combining the two sub-cases, we have
\[T_{3L}=\frac{q^2(q-1)^3}{2}+\frac{q^2(q-1)^2(q-2)^2}{6}.\]
Adding all contributions gives the stated formula.
\end{proof}

\end{document}
